% Part: applied-modal-logics
% Chapter: epistemic-logic
% Section: relational-models

\documentclass[../../../include/open-logic-section]{subfiles}

\begin{document}

\olfileid[tr]{aml}{el}{rel}

\olsection{Bağıntısal Modeller}

Epistemik mantıkların temel semantik kavramı, olağan modal mantıklarınkiyle
aynıdır. Bağıntısal modeller yine bir dünyalar kümesinden ve hangi
!!{propositional variable}s{}nin hangi dünyalarda ``doğru'' sayılacağını
belirleyen bir atamadan oluşur. Yalnızca tek bir özneyle ilgileniyorsak her
zamanki gibi tek bir erişilebilirlik bağıntımız vardır. Ancak çok özneli bir
epistemik mantığımız varsa bu tek erişilebilirlik bağıntısının yerini, her
$a$ için bir tane olmak üzere, özne simgeleri kümemiz $G$ ile dizinlenmiş bir
erişilebilirlik bağıntıları kümesi alır.

Bir \emph{bağıntısal model}, ikili erişilebilirlik bağıntılarıyla---her özne
için bir tane---birbirine bağlanan bir dünyalar kümesi ile hangi
!!{propositional variable}s{}nin hangi dünyalarda doğru olduğunu belirleyen bir
atamadan oluşur.

\begin{defn}
  Çok özneli epistemik dil için bir \emph{model}, $\mModel{M} = \tuple{W, R,
  V}$ üçlüsüdür; burada
  \begin{enumerate}
  \item $W$, boş olmayan bir ``dünyalar'' kümesidir,
  \item Her $a \in G$ için ${R}_a$, $W$ üzerinde ikili bir erişilebilirlik
  bağıntısıdır ve
  \item $V$, her !!{propositional variable}~$p$'ye bir $V(p)$ olası dünyalar
  kümesi atayan bir fonksiyondur.
  \end{enumerate}
  $R_a ww'$ geçerli olduğunda $w'$ dünyasının, a öznesi açısından,
  $w$'den \emph{erişilebilir} olduğunu söyleriz. $w \in V(p)$ olduğunda
  $p$'nin $w$'de \emph{doğru} olduğunu söyleriz.
  % [tr source emendation] Upstream's phrase `accessible by a from w` is
  % ungrammatical; Turkish states the intended indexed accessibility relation.
\end{defn}

İşleyiş, yalnızca daha fazla erişilebilirlik bağıntısı eklenmiş olması dışında
normal modal mantıktakiyle aynıdır. Belirli bir öznenin erişilebilirlik
bağıntısını genellikle onun bilgisel durumlarına ilişkin bir şeyi temsil
edecek biçimde yorumlarız. Örneğin çoğu kez $R_a ww'$ ifadesini, $w'$'nün
$a$'nın $w$'deki bilgisiyle bağdaşmasını ifade ediyor sayarız. Başka bir
deyişle özne $w$'de, $w$ dünyası ile $w'$ dünyasını birbirinden ayırt edemez.

\end{document}
